\documentclass[../main.tex]{subfiles}

\begin{document}

\section{LU Decomposition}

The last popular method to solve a linear system that we will
discuss is LU decomposition. As always, let's start by discussing
fundamental definitions.

\begin{definition}{}{}
\textbf{LU decomposition} is the process of factorizing an
invertible matrix into the product of lower and upper triangular
matrices where the diagonal elements of the lower triangular matrix
are $1$.
\end{definition}

Before we look at an important theorem and how it can be used to
solve a linear system, here is an example of LU decomposition
for matrix $A$ defined as the following.
\[
A =
\begin{bmatrix}
    1 & 2  & 3 \\
    2 & 6  & 9 \\
    3 & 10 & 18
\end{bmatrix}
\]
From here, notice that matrix $A$ can be transformed to upper
triangular matrix with the following elementary row operations of
$R_3$.
\[
\begin{bmatrix}
    1 & 2  & 3 \\
    2 & 6  & 9 \\
    3 & 10 & 18
\end{bmatrix}
\rightarrow
\begin{bmatrix}
    1 & 2  & 3 \\
    0 & 2  & 3 \\
    3 & 10 & 18
\end{bmatrix}
\rightarrow
\begin{bmatrix}
    1 & 2 & 3 \\
    0 & 2 & 3 \\
    0 & 4 & 9
\end{bmatrix}
\rightarrow
\begin{bmatrix}
    1 & 2 & 3 \\
    0 & 2 & 3 \\
    0 & 0 & 3
\end{bmatrix}
\]
Notice that the operations above can be represented as the
matrix multiplication with $R_3$ that are lower triangular matrices.
\[
\begin{bmatrix}
    1 & 0  & 0 \\
    0 & 1  & 0 \\
    0 & -2 & 1
\end{bmatrix}
\begin{bmatrix}
    1  & 0 & 0 \\
    0  & 1 & 0 \\
    -3 & 0 & 1
\end{bmatrix}
\begin{bmatrix}
    1  & 0 & 0 \\
    -2 & 1 & 0 \\
    0  & 0 & 1
\end{bmatrix}
A
=
\begin{bmatrix}
    1 & 2 & 3 \\
    0 & 2 & 3 \\
    0 & 0 & 3
\end{bmatrix}
\]
Therefore, the following equation is obtained.
\[
A
= \left(
    \begin{bmatrix}
        1 & 0 & 0 \\
        2 & 1 & 0 \\
        0 & 0 & 1
    \end{bmatrix}
    \begin{bmatrix}
        1 & 0 & 0 \\
        0 & 1 & 0 \\
        3 & 0 & 1
    \end{bmatrix}
    \begin{bmatrix}
        1 & 0 & 0 \\
        0 & 1 & 0 \\
        0 & 2 & 1
    \end{bmatrix}
\right)
\begin{bmatrix}
    1 & 2 & 3 \\
    0 & 2 & 3 \\
    0 & 0 & 3
\end{bmatrix}
=
\begin{bmatrix}
    1 & 0 & 0 \\
    2 & 1 & 0 \\
    3 & 2 & 1
\end{bmatrix}
\begin{bmatrix}
    1 & 2 & 3 \\
    0 & 2 & 3 \\
    0 & 0 & 3
\end{bmatrix}
\]
$A$ is now represented as the product of lower and upper triangular
matrices! Now that we know how LU decomposition works, let's take
a look at an important lemma on triangular matrices to prove a
theorem on LU decomposition.

\begin{lemma}{}{Inverse of Triangular Matrices}
The inverse of upper and lower triangular matrices are upper and
lower triangular matrices respectively.
\end{lemma}
\begin{proof}
First and foremost, consider an invertible upper triangular matrix
$A_{m \times m} = [a_{ij}]$ and its inverse $B = [b_{ij}]$.
By definition, its product $C = [c_{ij}]$ can be represented as
$C = [ \sum_{k = 1}^m a_{ij} b_{ji} ]$. Because $c_{ij} = 0$ for
$i \neq j$ and $c_{ij} = 1$ for $i = j$, the values of $b_{ji}$
can be found.

By definition, $a_{ij} = 0$ for $i > j$. Therefore, $c_{ij} = 0$
for $i > j$. Similarly, if $b_{ji} = 0$ for $j > i$, then $c_{ij} = 0$
for $i \neq j$. Continuing, if $b_{ij} = \frac{1}{a_{ij}}$ for
$i = j$, $B$ can be constructed, which is an upper triangular
matrix. Because the inverse matrix is unique, the inverse of an
upper triangular matrix is an upper triangular matrix if there exists
one.

Similarly, the fact that the inverse of an lower triangular matrix is
a lower triangular matrix can be proven. Let $A'_{m \times m} =
[a'_{ij}]$ be a lower triangular matrix and $B' = [b'_{ij}]$ be
its inverse. By definition, their product $C' = [c'_{ij}] =
[ \sum_{k = 1}^m a'_{ij} b'_{ji} ]$ is the identity matrix $I_m$.
Thus, it suffices to show that $c'_{ij} = 0$ for $i \neq j$ and
$c'_{ij} = 1$ for $i = j$.

Note that by definition, $a'_{ij} = 0$ for $i < j$. Therefore,
$c'_{ij} = 0$ for $i < j$. Similarly, if $b'_{ji} = 0$ for $i > j$,
and $b'_{ij} = \frac{1}{a'_{ij}}$ for $i = j$, then the valid
inverse matrix $B$ is constructed. Because $B$ is a lower triangular
matrix and the inverse matrix is unique, the inverse of a lower
triangular matrix is a lower triangular matrix.
\end{proof}

With the lemma in mind, let's prove the following theorem.

\begin{theorem}{}{Condition for LU Decomposition}
An invertible matrix has a LU decomposition if and only if
multiplying a sequence of lower triangular elementary matrices for
elementary row operations can transform the matrix into an upper
triangular matrix.
\end{theorem}
\begin{proof}
First, the statement that a nonsingular matrix has a LU decomposition
if multiplying a sequence of lower triangular elementary matrices for
elementary row operations can transform the matrix into an upper
triangular matrix can be proven. Let $A$ be an invertible matrix
that can be represented as the following.
\[
E_{n,n-1} \cdots E_{31} E_{21} A = U
\]
for lower matrices $E$ and an upper matrix $U$. Notice that $A$ can
be written as the following.
\[
A = \left( E^{-1}_{21} E^{-1}_{31} \cdots E^{-1}_{n,n-1} \right) U
\]
From Lemma \ref{lem:Inverse of Triangular Matrices}, the inverses of
the elementary matrices are lower triangular matrices. Moreover,
by Theorem \ref{theo:Property of Triangular Matrix Multiplication},
the product of the lower triangular matrices is also a lower
triangular matrix. Therefore, for a lower triangular matrix $L$,
$A$ can be written as the following.
\[
A = LU
\]
Therefore, the first half of the statement holds.

For the second half of the statement, consider the following equation.
\[
I = AA^{-1} = LUA^{-1} = L \left( UA^{-1} \right)
\]
Therefore, $L$ is nonsingular. From the proof for Lemma
\ref{lem:Inverse of Triangular Matrices}, it is evident that the diagonal
elements of an invertible lower triangular matrix are nonzero as
the diagonal elements in the inverse cannot be defined if a zero is
present in the diagonal elements of the lower triangular matrix.

Let $D$ be the diagonal matrix formed with the diagonal elements of
$L$ and define $L'$ such that the equation $L'D = L$ is satisfied.
Moreover, let $U' = DU$. Thus, the following equation is obtained.
\[
A = LU = L'DU = L'U'
\]
Recall that $L$ is invertible. Because $L^{-1}$ exists and
$(L'D)^{-1} = D^{-1} L'^{-1}$ exists, $L'$ is also invertible.
In other words, $L'^{-1} A = U'$. Therefore, it suffices to show
that $L'^{-1}$ is a product of lower triangular elementary matrices.
By definition of $D$, the diagonal elements of $L'$ are $1$. Thus,
it can be interpreted as a sequence of elementary row operations.
In other words, it can be represented as a product of lower triangular
elementary matrices, proving the latter half of the theorem.
\end{proof}

Now that we discussed what LU decomposition is, let's see how
we can use it to solve a linear system.

\subsection{Application in Linear Systems}

When we apply LU decomposition in solving linear systems, two
relatively easy systems must be solved. Consider the system
$Ax = b$ where $A = LU$. First, $y$ must be found from the
following system with forward substitution.
\[
Ly = b
\]
Then, we can solve for $x$ with back substitution with the following
system.
\[
Ux = y
\]
The good news is that both systems are relatively easy to solve.
The reason is simple. Consider the following equation.
\[
Ax = LUx = b
\]
Let $y = Ux$, then $Ly = b$. If we find $Ly = b$, then we can find
$Ux = y$ as we know $y$. This way, we can find $x$ while exploiting
the triangular form. Let's take a look at an example.

\begin{problem}{}{}
Solve the following system of linear equations.
\begin{align*}
    x - 2y + z &= 2 \\
    2x + y - z &= 1 \\
    3x - y + 2z &= 15
\end{align*}
\end{problem}
\begin{solution}
First and foremost, the linear system can be represented as the
matrix multiplication below.
\[
\begin{bmatrix}
    1 & -2 & 1 \\
    2 & 1  & -1 \\
    3 & -1 & 2
\end{bmatrix}
\begin{bmatrix}
    x \\
    y \\
    z
\end{bmatrix}
=
\begin{bmatrix}
    2 \\
    1 \\
    15
\end{bmatrix}
\]
Continuing, the coefficient matrix $A$ can be transformed as the
following.
\[
\begin{bmatrix}
    1 & -2 & 1 \\
    2 & 1  & -1 \\
    3 & -1 & 2
\end{bmatrix}
\rightarrow
\begin{bmatrix}
    1 & -2 & 1 \\
    0 & 5  & -3 \\
    3 & -1 & 2
\end{bmatrix}
\rightarrow
\begin{bmatrix}
    1 & -2 & 1 \\
    0 & 5  & -3 \\
    0 & 5  & -1
\end{bmatrix}
\rightarrow
\begin{bmatrix}
    1 & -2 & 1 \\
    0 & 5  & -3 \\
    0 & 0  & 2
\end{bmatrix}
\]
Using elementary matrices, the transformation above can be
written as the following.
\[
\begin{bmatrix}
    1 & 0  & 0 \\
    0 & 1  & 0 \\
    0 & -1 & 1
\end{bmatrix}
\begin{bmatrix}
    1  & 0 & 0 \\
    0  & 1 & 0 \\
    -3 & 0 & 1
\end{bmatrix}
\begin{bmatrix}
    1  & 0 & 0 \\
    -2 & 1 & 0 \\
    0  & 0 & 1
\end{bmatrix}
A =
\begin{bmatrix}
    1 & -2 & 1 \\
    0 & 5  & -3 \\
    0 & 0  & 2
\end{bmatrix}
\]
Completing the LU decomposition, $A$ can be represented as the
product of the following lower and upper triangular matrices.
\begin{align*}
    A
    &=
    \left(
    \begin{bmatrix}
        1 & 0 & 0 \\
        2 & 1 & 0 \\
        0 & 0 & 1
    \end{bmatrix}
    \begin{bmatrix}
        1 & 0 & 0 \\
        0 & 1 & 0 \\
        3 & 0 & 1
    \end{bmatrix}
    \begin{bmatrix}
        1 & 0 & 0 \\
        0 & 1 & 0 \\
        0 & 1 & 1
    \end{bmatrix}
    \right)
    \begin{bmatrix}
        1 & -2 & 1 \\
        0 & 5  & -3 \\
        0 & 0  & 2
    \end{bmatrix} \\
    &=
    \begin{bmatrix}
        1 & 0 & 0 \\
        2 & 1 & 0 \\
        3 & 1 & 1
    \end{bmatrix}
    \begin{bmatrix}
        1 & -2 & 1 \\
        0 & 5  & -3 \\
        0 & 0  & 2
    \end{bmatrix}
\end{align*}
Continuing, the following system could be solved.
\[
\begin{bmatrix}
    1 & 0 & 0 \\
    2 & 1 & 0 \\
    3 & 1 & 1
\end{bmatrix}
\begin{bmatrix}
    x' \\
    y' \\
    z'
\end{bmatrix}
=
\begin{bmatrix}
    2 \\
    1 \\
    15
\end{bmatrix}
\]
The matrix multiplication above can be written as the following
system of linear equations.
\begin{align*}
    x' &= 2 \\
    2x' + y' &= 1 \\
    3x' + y' + z' &= 15
\end{align*}
Substituting, $y' = 1 - 4 = -3$ and $z' = 15 - 6 + 3 = 12$. Thus,
the final system can be written as the following.
\[
\begin{bmatrix}
    1 & -2 & 1 \\
    0 & 5  & -3 \\
    0 & 0  & 2
\end{bmatrix}
\begin{bmatrix}
    x \\
    y \\
    z
\end{bmatrix}
=
\begin{bmatrix}
    2 \\
    -3 \\
    12
\end{bmatrix}
\]
Writing in linear system, $x$, $y$, and $z$ can be found.
\begin{align*}
    x - 2y + z &= 2 \\
    5y - 3z &= -3 \\
    2z &= 12
\end{align*}
Therefore, $z = 6$, $y = 3$, and $x = 2$.
\end{solution}

With LU decomposition, we completed our discussion on solving
systems of linear equations! Below are some practice problems
for the topics that we covered.

\end{document}


